\section{Introduction}
\label{sec:introduction}

The minimum weighted feedback arc set (MWFAS) problem asks for a minimum-weight set of arcs whose removal makes a nonnegative weighted directed graph acyclic, equivalently for a vertex ordering that minimizes the total weight of arcs directed backward with respect to that order. The problem is NP-hard in general \cite{K72} and arises across graph optimization, combinatorial ranking, and algorithm engineering wherever pairwise directed relations are cyclic. A directed graph may encode that task \(u\) should precede task \(v\), that a circuit signal must be routed before a dependent component is evaluated, that a package must be installed before another can be configured, or that one alternative is preferred to another in paired-comparison data. Rank aggregation from inconsistent pairwise information is a standard motivation for ordering heuristics \cite{ACN08}. In practice, directed relations are often cyclic, arising from noisy measurements, conflicting judgments, reciprocal dependencies, or competing design objectives. When arcs carry weights, backward constraints are not interchangeable: violating a high-weight precedence relation costs more than violating a weak one, and an effective heuristic must be designed around that asymmetry.

Weighted variants of MWFAS have an extensive heuristic and exact-method literature \cite{F90,BSNA21}. Exact methods are practical only for small instances or special structures. Naive score heuristics can be fast but ignore local cyclic structure; random multistart heuristics do not exploit graph decomposition; and specialized dense-ordering solvers may perform well on complete pairwise-comparison instances while being poorly aligned with sparse, irregular directed graphs. A meaningful structural distinction separates sparse general digraphs, where strongly connected components are scattered across a mostly acyclic background, from dense linear-ordering or tournament instances, where every ordered pair carries evidence and the problem is globally coupled. The computational question addressed here is whether a carefully engineered framework built on local-ratio seeding, topological add-back, and SCC-restricted refinement can deliver strong performance on nonnegative sparse weighted digraphs while clearly delineating its scope boundaries.

The local-ratio perspective is an important starting point for feedback arc set heuristics and approximation-oriented reasoning. We do not claim local-ratio as new. Broader local-ratio foundations appear in \cite{BYGR98}, and directed feedback problems were treated algorithmically in \cite{DF03}. The contribution in this paper is algorithm-engineering rather than a new approximation theorem: we study a reproducible local-ratio cycle-reduction implementation that records cycle-removal decisions, reconstructs an acyclic topological order, and then performs a topological add-back phase to recover arcs whenever doing so does not reintroduce cyclicity. This method is referred to as local-ratio topological add-back (LR-TA) after the components are defined in Section~\ref{sec:algorithmic-framework}. The emphasis is on deterministic data handling, measurable component effects, and integration with a refinement stage that is evaluated against exact and external baselines.

The integrated framework combines three components. First, the local-ratio topological add-back procedure produces a feasible ordering from weighted cycle reductions and a reconstruction step. Second, a WMSF-style weighted seed construction provides a complementary initial ordering and internal baseline. Third, incumbent-protected refinement searches within strongly connected components (SCCs). We refer to this refinement as incumbent-protected SCC neighborhood search (IPSNS). IPSNS updates the incumbent only when a candidate ordering has strictly smaller backward weight. This protection gives a simple invariant: the final solution cannot be worse than the best internal seed under the same objective and nonnegative-weight assumptions. The invariant is not an approximation guarantee, but it prevents refinement from degrading the seed and makes the local search behavior traceable instance by instance.

The computational study is organized around five evidence streams. The main sparse benchmark evaluates public weighted directed graph instances derived from graph-benchmark collections, with negative-weight instances separated from the standard nonnegative comparisons. An ablation study isolates the effect of the topological add-back phase, the use of the best seed, the incumbent-protected refinement, and the SCC-priority choice on a representative subset. An exact small-instance validation compares heuristic orderings against a bitmask dynamic-programming optimum where exact optimization is feasible. A sparse external-baseline study compares the proposed framework with the open-source DRMacIver/FAS heuristic, the Eades heuristic as implemented in python-igraph, a weighted adaptation of Eades et al. \cite{ELS93}, a Borda net-score ordering, and random multistart. Finally, a dense transfer study evaluates whether the framework carries over to complete weighted linear-ordering instances from the Linear Ordering Problem Library (LOLIB).

The results support a deliberately bounded claim. On 105 sparse benchmark instances, the incumbent-protected refinement has no incumbent violations relative to either internal seed. In the ablation study, the add-back phase reduces mean backward weight by 5.9 percent relative to local-ratio without add-back, and the refinement gives an additional 0.75 percent mean reduction on the tested subset. On the exact-validation subset of 57 standard nonnegative instances with \(n \le 20\), the refinement matches the dynamic-programming optimum on 56 instances, with mean relative gap about 0.0006 percent. On 97 standard nonnegative sparse instances in the external comparison, it achieves the best observed backward weight among the tested methods on 96 instances; the strongest external baseline, DRMacIver, is about 21.6 percent worse in mean backward weight on this sparse benchmark. These numbers support strong performance within the sparse nonnegative weighted-digraph setting studied here.

The dense LOLIB study is reported as a scope boundary rather than as primary evidence. LOLIB instances are complete weighted ordering instances, closely related to the linear ordering problem and structurally different from the sparse directed graphs that motivate the framework. On 50 dense LOLIB instances, the matrix-based pairwise-ordering baseline DRMacIver/FAS obtains the best observed solution on 45 instances, while the proposed refinement is best on 5 instances and has about 3.88 percent higher mean backward weight. This outcome is consistent with the structural distinction: SCC-local refinement is useful when sparse cyclic subgraphs provide meaningful neighborhoods, whereas dense complete ordering instances favor matrix-based pairwise-ordering heuristics. The method is therefore positioned as a reproducible general weighted-digraph framework with strong sparse performance and an explicitly documented dense-ordering limitation.

The contributions of this paper are as follows:
\begin{itemize}
    \item LR-TA is an engineered local-ratio seed with topological add-back: iterative cycle-weight reductions build a feasible ordering, and a heavy-first arc-recovery phase reduces unnecessary deletions.
    \item A complementary weighted seed provides an independent feasible ordering and an internal non-worsening baseline, ensuring the refinement begins from a strong alternative constructive solution.
    \item IPSNS is an incumbent-protected SCC-local destroy-and-repair refinement that concentrates search on strongly connected components with positive backward contribution and accepts a candidate ordering only when the global backward weight strictly decreases.
    \item The framework provides correctness guarantees: every returned ordering is acyclic, and the final solution is never worse than the best internal seed under the same nonnegative-weight objective.
    \item An extensive computational validation covers exact bitmask dynamic-programming certification on small instances, time-capped MIP validation on medium instances, five external and adapted baselines, a component ablation study, a runtime--budget study, and a dense LOLIB transfer test.
    \item A fully reproducible artifact accompanies the paper, including all benchmark instances, algorithm implementations, experimental scripts, and saved result summaries.
\end{itemize}

The remainder of the paper is organized as follows. Section~\ref{sec:related-work} reviews feedback arc set heuristics, local-ratio methods, ordering heuristics, exact methods, and linear-ordering benchmarks. Section~\ref{sec:problem-definition} defines the minimum weighted feedback arc set objective, the ordering-induced backward weight, and the relation to dense linear-ordering notation. Section~\ref{sec:algorithmic-framework} describes the local-ratio topological add-back procedure, the weighted seed, and the incumbent-protected SCC refinement. Section~\ref{sec:experimental-design} presents datasets, baselines, metrics, and reproducibility controls. Section~\ref{sec:results} summarizes the computational results. Section~\ref{sec:discussion} discusses limitations, including the dense LOLIB boundary, and Section~\ref{sec:conclusion} concludes.
